\documentclass[11pt]{article}
\usepackage[margin=0.82in]{geometry}
\usepackage{amsmath,amssymb,amsthm,mathtools,microtype}
\usepackage[T1]{fontenc}
\usepackage{lmodern}
\usepackage[colorlinks=true,urlcolor=blue,linkcolor=blue]{hyperref}
\newtheorem{theorem}{Theorem}
\newcommand{\Z}{\mathbb Z}
\newcommand{\tr}{\operatorname{tr}}
\newcommand{\diag}{\operatorname{diag}}
\title{A braid-independent obstruction for the Ren core}
\author{Research note from the Jones unknot-detection investigation}
\date{October 31, 2026}
\begin{document}
\maketitle
\noindent\textbf{Result and scope.}
For one explicitly specified four-strand tangle, \emph{every} attached four-strand braid has annular bracket degree exactly two. Thus arbitrary braiding cannot cancel its last nonconstant annular term. This is a theorem about a particular infinite construction family. It neither proves Jones unknot detection nor exhibits a counterexample. No claim of priority in the literature is made.

\medskip
Let $\delta=-A^2-A^{-2}$, use the crossing convention
$\sigma_i=A\,1+A^{-1}e_i$, and normalize a contractible circle to $\delta$.
In the six-strand braid
\[
 b=\sigma_2^{-1}\sigma_1^3\sigma_2^{-1}
       \sigma_4\sigma_5^{-3}\sigma_4,
\]
cap positions $3,4$ together at the top and at the bottom, leaving the four outside endpoints at each end. Denote the resulting $(4,4)$-tangle by $T$. This is the core used in the construction of Ren \cite{Ren}; the two additional crossings for that example correspond to the attached braid $\sigma_1\sigma_3$. Write $\widehat{T\beta}_{\rm ann}$ for annular closure after attaching $\beta\in B_4$. The variable $z$ is an essential circle in the annulus.

\begin{theorem}
For every $\beta\in B_4$ there exist $a_\beta,b_\beta\in\Z[A^{\pm1}]$ such that
\[
 \big\langle\widehat{T\beta}_{\rm ann}\big\rangle
       =a_\beta(A)z^2+b_\beta(A),\qquad a_\beta(A)\ne0.
\]
This includes links of any number of components and braids of arbitrary length.
\end{theorem}

\begin{proof}
The coefficient of the four-through identity diagram in $T\in TL_4$ is zero. This follows by expanding the ten crossings above and capping, and is also the highest-degree cancellation underlying Ren's example. The other thirteen diagrams span a two-sided ideal, so $T\beta$ has no four-through component.

The two-through standard module of $TL_4$ has dimension three. Its generator matrices are
\[
 E_1=\begin{pmatrix}\delta&1&0\\0&0&0\\0&0&0\end{pmatrix},\quad
 E_2=\begin{pmatrix}0&0&0\\1&\delta&1\\0&0&0\end{pmatrix},\quad
 E_3=\begin{pmatrix}0&0&0\\0&0&0\\0&1&\delta\end{pmatrix}.
\]
Denote this representation by $\rho$. On each of the thirteen nonidentity diagrams, its matrix trace equals the coefficient of $z^2$ in annular closure. Equivalently, this is the two-through character in the annular trace decomposition. Hence
\[
 a_\beta=\tr\bigl(\rho(T)\rho(\beta)\bigr).
\]
\newpage
Put $x=-A^{-4}$ and $S=\diag(1,A^2,A^4)$. After removing each crossing's scalar $A$, the conjugated generator matrices become
\[
 R_1=\begin{pmatrix}x&1&0\\0&1&0\\0&0&1\end{pmatrix},\quad
 R_2=\begin{pmatrix}1&0&0\\-x&x&1\\0&0&1\end{pmatrix},\quad
 R_3=\begin{pmatrix}1&0&0\\0&1&0\\0&-x&x\end{pmatrix}.
\]
They and their inverses have entries in $\Z[x^{\pm1}]$, since their determinants are $x$. If $e(\beta)$ is the braid exponent sum, define
\[
 U(x)=A^{-e(\beta)}S^{-1}\rho(\beta)S,\qquad
 M(x)=A^{-2}S^{-1}\rho(T)S.
\]
Both are Laurent matrices over $\Z$. In particular,
\[
 a_\beta=A^{e(\beta)+2}f(x),\qquad f(x)=\tr(M(x)U(x)).
\]
The explicit core calculation (entries are listed on the next page) gives
\[
 M_0:=M(1)=
 \begin{pmatrix}-15&0&15\\-25&0&25\\-15&0&15\end{pmatrix},\qquad
 M_1:=M'(1)=
 \begin{pmatrix}-33&0&0\\-40&-8&-15\\-15&0&-18\end{pmatrix}.
\]
Thus $M_0\equiv0\pmod5$ and $M_1\equiv2I\pmod5$.

Set $u=(1,0,1)^t$, $v=(1,0,-1)$, and $e_2=(0,1,0)^t$.
Every $R_i(1)$ has determinant one and fixes $u$. The same is therefore true of $U_0:=U(1)\in SL_3(\Z)$. Since
\[
 M_0=(-15u-25e_2)v,
 \qquad f(1)=-25\,vU_0e_2,
\]
the condition $f(1)=0$ implies $vU_0e_2=0$.

The primitive vector $u$ generates a saturated copy of $\Z$ in $\Z^3$.
The quotient $\Z^3/\Z u$ has basis $[e_1],[e_2]$, and $v$ gives its first-coordinate functional. The induced matrix of $U_0$ is integral with determinant one. If $vU_0e_2=0$, its upper-right entry is zero, so its two diagonal entries are either both $1$ or both $-1$. Consequently
\[
 \tr U_0\in\{3,-1\}.
\]
The derivative $U'(1)$ is integral. Differentiating $f$ now gives
\[
 f'(1)=\tr(M_1U_0+M_0U'(1))
      \equiv 2\tr U_0
      \equiv1\ \text{or}\ 3\pmod5.
\]
Thus $f(1)$ and $f'(1)$ cannot both vanish. In particular $f$ is not the zero Laurent polynomial, proving $a_\beta\ne0$.
\end{proof}

\noindent\textbf{Stronger local conclusion.}
The normalized coefficient $f(x)$ has order of vanishing at most one at $x=1$. This is a universal obstruction; it does not infer a polynomial identity from finitely many matching Taylor coefficients.

\newpage
\noindent\textbf{Explicit finite certificate.}
With the conventions above, the only nonzero entries of $M(x)$ are
\begin{align*}
 M_{11}&=-2-3x-4x^2-3x^3-2x^4-x^5,\\
 M_{13}&=x^{-3}+2x^{-2}+3x^{-1}+3+3x+2x^2+x^3,\\
 M_{21}&=-x^{-2}-2x^{-1}-4-5x-5x^2-4x^3-3x^4-x^5,\\
 M_{22}&=x^{-2}+x^{-1}-x^2-x^3,\\
 M_{23}&=x^{-4}+3x^{-3}+4x^{-2}+5x^{-1}+5+4x+2x^2+x^3,\\
 M_{31}&=-x^{-2}-2x^{-1}-3-3x-3x^2-2x^3-x^4,\\
 M_{33}&=x^{-4}+2x^{-3}+3x^{-2}+4x^{-1}+3+2x.
\end{align*}
The two omitted entries $M_{12},M_{32}$ are zero. Substitution and differentiation directly reproduce $M_0,M_1$. The accompanying exact-arithmetic verifier rebuilds $T$ from the stated ten-crossing word, checks the absent identity coefficient, constructs the module, and verifies its trace against annular closure on all thirteen relevant basis diagrams.

\medskip
\noindent\textbf{Independent checks.}
Two separate algebraic audits checked the change of basis, trace interpretation, primitive quotient lattice, and arithmetic argument. Before the proof was found, an exhaustive modular computation checked all $366{,}210{,}937$ freely reduced words of length at most twelve. Among them were $183{,}771{,}832$ instances whose closures are knots; none had vanishing $z^2$ coefficient at $A=2$ modulo $1{,}000{,}000{,}007$. These are word and closure-instance counts, not counts of distinct knots. The search implementation was compared with exact Laurent-polynomial calculations on 128 independent sample words. The theorem above does not depend on this bounded search.

\medskip
\noindent\textbf{What this does and does not eliminate.}
The result rules out obtaining an annular-degree-zero pattern by modifying only the attached four-strand braid of this core. It does not imply that every ordinary closure has nontrivial Jones polynomial: filling the annulus substitutes $z=\delta$, and cancellation after that substitution is a different equation. Indeed, an ordinary-unknot closure already occurs in this family.

The earlier satellite condition involving the normalized three-dimensional colored Jones polynomial $H_3(C)=1$ remains unresolved. It can also be converted to ordinary Jones polynomial one by a $(2,1)$-cable, so the present core is not essential to that route. None of the calculations in this pass supplied a nontrivial companion satisfying $H_3(C)=1$.

\begin{thebibliography}{9}
\bibitem{Ren} Q. Ren, \emph{A counterexample to the wrapping number conjecture},
arXiv:2609.11084v2, September 14, 2026.
\url{https://arxiv.org/html/2609.11084v2}.
\end{thebibliography}
\end{document}